% ===========================================================================
\section{The four subspaces, in one breath}\label{sec:subspaces}
\warmth{0}\quad\whisper{$U$ splits the output space, $V$ splits the input space, along the rank.}

\begin{strand}
Cut the singular vectors at the rank $r$ (the number of nonzero $\sigma_i$). The first $r$
columns of $V$ and the last $n-r$ split the input $\R^n$ into two orthogonal pieces; the first
$r$ columns of $U$ and the last $m-r$ split the output $\R^m$. Those four pieces are exactly the
\keyword{four fundamental subspaces} of $A$, and the SVD names all of them at once.
\end{strand}

\begin{theorem}[the four subspaces]\label{thm:foursub}
Let $A=U\Sig V\T$ have rank $r$ (so $\sigma_1,\dots,\sigma_r>0$ and the rest vanish). Then:
\end{theorem}
\noindent\begin{tabularx}{\linewidth}{@{}l l L@{}}
\toprule
{\color{indigo}subspace} & {\color{indigo}lives in} & \multicolumn{1}{l}{\color{indigo}is the span of \dots \;(dimension)}\\
\midrule
column space $\col(A)$        & $\R^m$ & $u_1,\dots,u_r$ \;($\dim=\fillin[0.8cm]$)\\
left null space $\Null(A\T)$  & $\R^m$ & $u_{r+1},\dots,u_m$ \;($\dim=m-r$)\\
row space $\col(A\T)$         & $\R^n$ & $\fillin[2.5cm]$ \;($\dim=r$)\\
null space $\Null(A)$         & $\R^n$ & $v_{r+1},\dots,v_n$ \;($\dim=\fillin[1.2cm]$)\\
\bottomrule
\end{tabularx}

\begin{strand}
The picture behind the table: $A$ sends the row space \emph{isomorphically} onto the column
space, stretching the $i$-th axis by $\sigma_i$, and it \emph{annihilates} the null space. Nothing
else happens. Rank--nullity ($\dim\Null(A)=n-r$) and the orthogonal splittings
$\col(A)\perp\Null(A\T)$ in $\R^m$, $\col(A\T)\perp\Null(A)$ in $\R^n$ are now visibly true,
because the columns of an orthogonal matrix are orthonormal.
\end{strand}

\recall{Rank, finally honest: $\rank(A)=$ the number of nonzero singular values. Tiny-but-nonzero $\sigma_i$ are why ``numerical rank'' is a threshold on $\Sig$, not a yes/no.}

\begin{yourturn}
Take $\displaystyle A=\begin{pmatrix}1&1\\1&1\end{pmatrix}$, with SVD given by $\sigma_1=2,\ \sigma_2=0$,
$\,u_1=v_1=\tfrac1{\sqrt2}(1,1)$, $\,v_2=\tfrac1{\sqrt2}(1,-1)$. Read off:
\[
  \rank(A)=\fillin[0.8cm],\quad \col(A)=\operatorname{span}\fillin[2cm],\quad
  \Null(A)=\operatorname{span}\fillin[2cm].
\]
\recall{Rank $1$; column space $=\operatorname{span}\{(1,1)\}$; null space $=\operatorname{span}\{(1,-1)\}$. The map flattens the plane onto a line.}
\end{yourturn}

\begin{strand}
Pseudoinverse preview: because $A$ is a clean isomorphism row space $\to$ column space, it has an
obvious inverse \emph{there} (divide by $\sigma_i$) and zero elsewhere. That is the Moore--Penrose
pseudoinverse of §\ref{sec:faces}.
\end{strand}
